
\documentclass[journal,12pt,onecolumn]{IEEEtran}
\IEEEoverridecommandlockouts
\usepackage{cite}

\usepackage{amsmath,amssymb,amsfonts}
\usepackage{algorithmic}
\usepackage{algorithm}
%\usepackage{algpseudocode}
\usepackage{graphicx}

\usepackage{booktabs} 

\usepackage{textcomp}
%\usepackage[marginal]{footmisc}
\usepackage{xcolor}
\usepackage{subfigure}
\graphicspath{{./img/}}
\usepackage{tabularx}
\usepackage{makecell}
\usepackage{caption2}
%\usepackage{fontspec}
%\setmainfont{TeX Gyre Termes}
\usepackage{color}

\usepackage[UTF8]{ctex}


\def\BibTeX{{\rm B\kern-.05em{\sc i\kern-.025em b}\kern-.08em
		T\kern-.1667em\lower.7ex\hbox{E}\kern-.125emX}}
\begin{document}
\title{SNN Information Reconstruction}
\author{\IEEEauthorblockN{author}
\\
	\IEEEauthorblockA{IEEEauthorblockA}
}

\maketitle

\begin{abstract}
  abstract abstract abstract
\\ 
\end{abstract}
\begin{IEEEkeywords}


  IEEEkeywords IEEEkeywords IEEEkeywords
	
\end{IEEEkeywords}

\let\thefootnote\relax


\section{Introduction}
Introduction 

\cite{apolinario2025tess}



\section{时间反向传播（BPTT）：完整时序梯度计算}

BPTT（Backpropagation Through Time）是训练时序神经网络（RNN/SNN）的标准方法，其核心思想是：
\textbf{通过时间展开，将递归网络转化为深度前馈网络，并计算完整的时空梯度。}

\subsection{网络动态（前向过程）}

考虑一个时间序列输入 $\{x[t]\}_{t=1}^T$，神经元状态满足：

\[
h[t] = f(h[t-1], x[t]; \theta)
\]

输出为：

\[
y[t] = g(h[t]; \theta)
\]

总损失为：

\[
L = \sum_{t=1}^{T} \ell(y[t], y^*[t])
\]

\subsection{关键问题：时间信用分配（Temporal Credit Assignment）}

由于当前状态依赖历史：

\[
h[t] \rightarrow h[t+1] \rightarrow \cdots \rightarrow h[T]
\]

因此权重 $W$ 在时间 $s$ 的影响会传递到未来所有时间：

\[
\frac{\partial L}{\partial W}
= \sum_{t=1}^{T} \frac{\partial L}{\partial h[t]} \frac{\partial h[t]}{\partial W}
\]

\subsection{完整梯度展开（核心公式）}

展开时间依赖：

\[
\frac{\partial L}{\partial W}
=
\sum_{t=1}^{T}
\frac{\partial L}{\partial h[t]}
\left(
\sum_{s=1}^{t}
\frac{\partial h[t]}{\partial h[s]}
\frac{\partial h[s]}{\partial W}
\right)
\]

含义：

- 内层求和：时间依赖传播（history effect）

- 外层求和：所有时间步的损失贡献

\subsection{反向传播递推（核心递推公式）}

定义误差信号：

\[
\delta[t] = \frac{\partial L}{\partial h[t]}
\]

递推关系：

\[
\delta[t] =
\frac{\partial L[t]}{\partial h[t]}
+
\delta[t+1]
\frac{\partial h[t+1]}{\partial h[t]}
\]

说明：

- 梯度从 $T \rightarrow 1$ 反向传播
- 每一步累积未来影响

\subsection{参数梯度}

\[
\frac{\partial L}{\partial W}
=
\sum_{t=1}^{T}
\delta[t]
\frac{\partial h[t]}{\partial W}
\]


\subsection{问题与局限}

- 需要存储所有时间状态（内存 $O(T)$）

- 无法在线学习

- 梯度消失/爆炸

BPTT 能够精确地计算包含长期依赖的梯度，但在长序列上存在梯度消失和爆炸的问题。梯度消失会导致早期时间步贡献极小，而梯度爆炸会导致训练不稳定。

\section{S‑TLLR方法的完整推导与分析\cite{apolinario2023s}}

\subsection{神经元动力学模型}

我们使用离散时间下的LIF（泄露积分发放）神经元模型，该模型在论文中作为基础动力学定义。脉冲神经元状态更新公式为：

\begin{equation}
u_i[t] = \gamma (u_i[t-1] - v_{th} \, y_i[t-1]) + \sum_j w_{ij}\, x_j[t],
\end{equation}

\begin{equation}
y_i[t] = H(u_i[t] - v_{th}),
\end{equation}

其中：
\begin{itemize}
    \item $u_i[t]$ 表示第 $i$ 个神经元在时间 $t$ 的膜电位；
    \item $y_i[t]$ 是基于膜电位的二值脉冲输出；
    \item $\gamma$ 是膜电位的泄露系数；
    \item $H(\cdot)$ 是阶跃函数；
    \item $w_{ij}$ 是从第 $j$ 个前馈神经元到第 $i$ 个神经元的权重；
    \item $x_j[t]$ 是前一层的输入脉冲。
\end{itemize}

这一模型表明：膜电位既受到自身历史状态累积的影响，也受到当前时刻输入的驱动，这样的时序依赖是SNN训练难点的基础。

\subsection{三因子学习规则的形式}

S‑TLLR属于广义三因子学习规则，即权重更新由三个信号共同决定：

\begin{equation}
\Delta w_{ij}[t] = \delta_i[t] \cdot e_{ij}[t],
\end{equation}

其中：
\begin{itemize}
    \item $\delta_i[t]$ 是自顶向下的学习信号（第三因子），负责将全局监督误差信息传递至局部；
    \item $e_{ij}[t]$ 是资格迹（eligibility trace），表示前后脉冲序列之间的时间相关性；
    \item 这一乘积形式使得学习同时考虑误差信号和局部时序痕迹。
\end{itemize}

三因子结构是许多生物启发学习机制的基础，其优点是局部可计算、可用于在线学习。

\subsection{STDP启发的资格迹定义}

S‑TLLR 引入了受 STDP 原理启发的资格迹表达式，既包含“因果”时间相关成分，也包含“非因果”成分：

\begin{equation}
e_{ij}[t] =
\alpha_{\text{pre}} \sum_{\tau=0}^{t} \lambda^{\,t-\tau} \, x_j[\tau] \, \Psi(u_i[t])
+
\alpha_{\text{post}} \sum_{\tau=0}^{t-1} \lambda^{\,t-\tau} \, \Psi(u_i[\tau]) \, x_j[t].
\label{eq:stllr_eligibility}
\end{equation}

其中：
\begin{itemize}
    \item $e_{ij}[t]$ 表示从神经元 $j$ 到神经元 $i$ 在时间 $t$ 的资格迹；
    \item $\lambda$ 是指数衰减因子（类似于时间常数）；
    \item $\alpha_{\text{pre}}$ 和 $\alpha_{\text{post}}$ 分别是预反馈和后反馈的权重参数；
    \item $\Psi(u)$ 是次级激活函数（surrogate function），用于平滑处理 $u$ 对学习规则的影响；
    \item 第一项代表因果贡献：过去的前馈输入 $x_j[\tau]$ 影响当前膜电位；
    \item 第二项代表非因果贡献：当前前馈 spike 与过去后膜电位的相关性。
\end{itemize}

该资格迹将 STDP 的时间先后关系转化为可累积的因果/非因果成分，从而对时序信息进行编码。

\subsection{公式中各部分的深层含义}

对于第一项：

\begin{equation}
\alpha_{\text{pre}} \sum_{\tau=0}^{t} \lambda^{\,t-\tau} \, x_j[\tau] \, \Psi(u_i[t]),
\end{equation}

它捕捉了历史前馈输入 $x_j[\tau]$ 随时间衰减对当前膜电位的影响，这是 STDP 中“pre leads post” 的生物现象。

第二项：

\begin{equation}
\alpha_{\text{post}} \sum_{\tau=0}^{t-1} \lambda^{\,t-\tau} \, \Psi(u_i[\tau]) \, x_j[t],
\end{equation}

则将当前的前馈 spike 与过去的后膜电位相乘，这是一种“非因果”关联，其重要性在于增强时间相关性信息的捕获。

这两部分共同构成了S‑TLLR的资格迹，其优势在于：

1. \textbf{保留了前后脉冲时间相关性}（比纯因果 STDP 更强）；  
2. \textbf{无需展开时间，可以在线计算}；  
3. \textbf{比传统三因子规则更简洁、内存复杂度低至 $O(n)$。} 

\subsection{学习信号的来源}

学习信号 $\delta_i[t]$ 有两种生成方式：

\begin{itemize}
    \item 利用误差反向传播产生的梯度信号（基于 BPTT 或逐层 BP）；
    \item 利用固定随机反馈矩阵直接从输出层生成学习信号（Random Feedback Alignment）。
\end{itemize}

这使得 S‑TLLR \textbf{既可以在监督学习场景中使用反向误差信号}，也可以在更生物可实现的随机反馈场景中估计全局误差信息。

\subsection{完整更新算法公式}

在时间循环中，S‑TLLR 的权重更新可以写为：

\begin{equation}
\Delta w_{ij}^{(l)}[t] = \delta_i^{(l)}[t] \, e_{ij}^{(l-1)}[t],
\end{equation}

该公式由两部分组成：
\begin{itemize}
    \item $\delta_i^{(l)}[t]$ 是来自上一层的学习信号，
    \item $e_{ij}^{(l-1)}[t]$ 是资格迹（由式\eqref{eq:stllr_eligibility}定义）。
\end{itemize}

在第 $L$ 层（输出层）中更新则为：

\begin{equation}
\Delta w_{ij}^{(L)}[t] = \delta_i^{(L)}[t] \, y_j^{(L-1)}[t],
\end{equation}

其中 $y_j^{(L-1)}[t]$ 是上一层的输出 spike。

在隐藏层：

\begin{equation}
\Delta w_{ij}^{(l)}[t] = \delta_i^{(l)}[t] \, e_{ij}^{(l-1)}[t],
\end{equation}

这一结构体现出 \textbf{三因子更新依赖局部时序信息和全局误差信号}。

\subsection{递推计算与内存复杂度}

虽然资格迹中包含历史求和，但可以用递推方式实现在线更新：

\begin{equation}
\text{(pre causal trace)}\quad a_{j}^{\text{pre}}[t] = \lambda a_{j}^{\text{pre}}[t-1] + x_j[t],
\end{equation}

\begin{equation}
\text{(post non‑causal trace)}\quad b_{i}^{\text{post}}[t] = \lambda b_{i}^{\text{post}}[t-1] + \Psi(u_i[t]),
\end{equation}

再组合：

\begin{equation}
e_{ij}[t] = \alpha_{\text{pre}}\,a_{j}^{\text{pre}}[t]\,\Psi(u_i[t])
+
\alpha_{\text{post}}\,b_{i}^{\text{post}}[t-1]\,x_j[t].
\end{equation}

这样实现表明 S‑TLLR 所需内存仅与神经元数线性相关，而不是与时间步成比例。

\subsection{方法的整体作用与优势}

S‑TLLR 的设计优势在于：

\begin{itemize}
    \item \textbf{时间 locality}：资格迹只依赖过去活动，不依赖未来；
    \item \textbf{空间 locality}：权重更新只需要局部 synapse 信息，不依赖全局梯度展开；
    \item \textbf{非因果整合}：同时考虑了因果和非因果的 STDP 型信息，有助于更稳定的学习；
    \item \textbf{低内存复杂度}：资格迹可在线递推计算，不需要保存整段历史。
\end{itemize}

这些特性使 S‑TLLR 适合在低资源、在线实时场景中训练深度 SNN，而不需要昂贵的 BPTT。


\section{S-TLLR与TESS方法的统一分析与对比}

\subsection{问题背景：时空信用分配}

在脉冲神经网络中，梯度可写为：

\begin{equation}
\frac{\partial L}{\partial W}
=
\sum_{\tau}
\underbrace{s[\tau]}_{\text{历史输入}}
\underbrace{\sum_{t \ge \tau} \delta[t] \cdot \frac{\partial u[t]}{\partial u[\tau]}}_{\text{时间信用分配}}
\end{equation}

其中：
\begin{itemize}
\item 时间维度：$\frac{\partial u[t]}{\partial u[\tau]}$
\item 空间维度：$\delta[t]$ 来自其他层
\end{itemize}

因此学习问题本质是：
\begin{equation}
\text{Credit Assignment} = \text{Temporal} + \text{Spatial}
\end{equation}



\subsection{S-TLLR的推导与结构}

\subsubsection{三因子学习规则}

S-TLLR采用三因子学习框架：

\begin{equation}
\Delta W_{ij}[t] = \delta_i[t] \cdot e_{ij}[t]
\end{equation}

其中：
\begin{itemize}
\item $e_{ij}[t]$：eligibility trace（局部时间信息）
\item $\delta_i[t]$：误差信号（非局部）
\end{itemize}



\subsubsection{Eligibility Trace的定义}

S-TLLR构造如下资格迹：

\begin{equation}
e_{ij}[t] =
\alpha_{\text{pre}} \sum_{\tau \le t} \lambda^{t-\tau} x_j[\tau] \Psi(u_i[t])
+
\alpha_{\text{post}} \sum_{\tau < t} \lambda^{t-\tau} \Psi(u_i[\tau]) x_j[t]
\end{equation}

该表达式可拆为：

\begin{equation}
e_{ij}[t] = e^{\text{causal}} + e^{\text{non-causal}}
\end{equation}



\subsubsection{递推形式}

定义：

\begin{equation}
a_j[t] = \lambda a_j[t-1] + x_j[t]
\end{equation}

\begin{equation}
b_i[t] = \lambda b_i[t-1] + \Psi(u_i[t])
\end{equation}

则：

\begin{equation}
e_{ij}[t] =
\alpha_{\text{pre}} a_j[t] \Psi(u_i[t])
+
\alpha_{\text{post}} b_i[t-1] x_j[t]
\end{equation}



\subsubsection{S-TLLR的关键特点}

\begin{itemize}
\item 时间局部性：
\begin{equation}
e_{ij}[t] \text{ 仅依赖 } t \text{ 及过去}
\end{equation}

\item 空间非局部性：
\begin{equation}
\delta_i[t] = \text{通过反向传播获得}
\end{equation}
\end{itemize}

因此：

\begin{equation}
\text{S-TLLR} = \text{Temporal Local} + \text{Spatial Non-local}
\end{equation}



\subsection{TESS的推导与结构}

TESS同样基于三因子学习：

\begin{equation}
\Delta W_{ij}[t] = m_i[t] \cdot e_{ij}[t]
\end{equation}

关键区别在于学习信号 $m_i[t]$。



\subsubsection{Temporal部分（与S-TLLR相同）}

TESS直接继承S-TLLR的资格迹：

\begin{equation}
e_{ij}[t] =
\alpha_{\text{pre}} a_j[t] \Psi(u_i[t])
+
\alpha_{\text{post}} b_i[t-1] x_j[t]
\end{equation}

即：

\begin{equation}
\text{Temporal part (TESS)} = \text{S-TLLR}
\end{equation}



\subsubsection{Spatial部分（核心创新）}

TESS不使用反向传播，而是构造：

\begin{equation}
m^{(l)}[t] = B^{(l)\top} \left( f(B^{(l)} o^{(l)}[t]) - y^* \right)
\end{equation}

其中：
\begin{itemize}
\item $B^{(l)}$：固定随机投影矩阵
\item $f(\cdot)$：任务函数（如softmax）
\item $y^*$：标签
\end{itemize}



\subsubsection{推导解释}

定义中间变量：

\begin{equation}
z^{(l)}[t] = B^{(l)} o^{(l)}[t]
\end{equation}

则误差：

\begin{equation}
\epsilon^{(l)}[t] = f(z^{(l)}[t]) - y^*
\end{equation}

再投影回神经元空间：

\begin{equation}
m^{(l)}[t] = B^{(l)\top} \epsilon^{(l)}[t]
\end{equation}



\subsubsection{关键性质}

该构造满足：

\begin{equation}
m^{(l)}[t] \approx \delta^{(l)}[t]
\end{equation}

但完全不依赖：

\begin{equation}
\frac{\partial L}{\partial u^{(l)}}
\text{ 的反向传播}
\end{equation}



\subsubsection{TESS的结构总结}

\begin{equation}
\text{TESS} = \text{Temporal Local} + \text{Spatial Local}
\end{equation}



\subsection{两者的统一表达}

将两者统一写为：

\begin{equation}
\Delta W_{ij}[t] = \underbrace{e_{ij}[t]}_{\text{历史信息}} \cdot \underbrace{\eta_i[t]}_{\text{误差信号}}
\end{equation}

其中：

\begin{equation}
\eta_i[t] =
\begin{cases}
\delta_i[t] & \text{S-TLLR} \\
m_i[t] & \text{TESS}
\end{cases}
\end{equation}



\subsection{信息流对比（核心结论）}

\subsubsection{S-TLLR}

\begin{equation}
\eta_i[t] = \delta_i[t]
\end{equation}

依赖：

\begin{equation}
\delta_i[t] =
\sum_k W_{ki}^{(l+1)} \delta_k^{(l+1)}[t]
\end{equation}

特点：

\begin{itemize}
\item 需要跨层传播
\item 非局部
\end{itemize}



\subsubsection{TESS}

\begin{equation}
\eta_i[t] = m_i[t]
\end{equation}

仅依赖：

\begin{equation}
o^{(l)}[t], \quad y^*
\end{equation}

特点：

\begin{itemize}
\item 完全局部生成
\item 无跨层依赖
\end{itemize}



\subsection{复杂度对比}

\begin{equation}
\text{Mem}_{\text{S-TLLR}} = O(Ln)
\end{equation}

\begin{equation}
\text{Mem}_{\text{TESS}} = O(Ln)
\end{equation}

但计算复杂度：

\begin{equation}
\text{Time}_{\text{S-TLLR}} = O(Ln^2)
\end{equation}

\begin{equation}
\text{Time}_{\text{TESS}} = O(LCn)
\end{equation}

其中 $C \ll n$。



\subsection{方法本质总结}

\begin{itemize}
\item S-TLLR：
\begin{equation}
\text{使用真实误差} + \text{局部时间信息}
\end{equation}

\item TESS：
\begin{equation}
\text{使用局部近似误差} + \text{局部时间信息}
\end{equation}
\end{itemize}



\subsection{最终核心差异}

\begin{equation}
\text{S-TLLR} \Rightarrow \text{保留空间梯度信息}
\end{equation}

\begin{equation}
\text{TESS} \Rightarrow \text{近似空间梯度信息}
\end{equation}

\section{NDOT方法的完整推导与机制分析}

\subsection{问题背景：BPTT的梯度结构}

对于多层SNN，BPTT的梯度可写为：

\begin{equation}
\frac{\partial L}{\partial W^l}
=
\sum_{t=1}^{T}
\frac{\partial L}{\partial s^l[t]}
\frac{\partial s^l[t]}{\partial u^l[t]}
\left(
\frac{\partial u^l[t]}{\partial W^l}
+
\sum_{k<t}
\prod_{i=k}^{t-1}
\left(
\frac{\partial u^l[i+1]}{\partial u^l[i]}
+
\frac{\partial u^l[i+1]}{\partial s^l[i]}
\frac{\partial s^l[i]}{\partial u^l[i]}
\right)
\frac{\partial u^l[k]}{\partial W^l}
\right)
\end{equation}



\subsubsection{含义分析}

该公式包含两个关键部分：

\begin{itemize}
\item 空间梯度：
\[
\frac{\partial L}{\partial s^l[t]}
\frac{\partial s^l[t]}{\partial u^l[t]}
\]

\item 时间梯度：
\[
\prod_{i=k}^{t-1} (\cdot)
\]
\end{itemize}



\subsubsection{问题本质}

\begin{equation}
\text{BPTT困难点：}  \text{时间展开、}  \text{存储所有历史状态}
\end{equation}



\subsection{NDOT核心思想：梯度分解}

NDOT提出：

\begin{equation}
\frac{\partial L}{\partial W^l}
=
\sum_{t=1}^{T}
g_u^l[t] \cdot \hat{a}^{l-1}[t]^T
\end{equation}

其中：

\begin{equation}
g_u^l[t] =
\left(
\frac{\partial L}{\partial s^l[t]}
\frac{\partial s^l[t]}{\partial u^l[t]}
\right)^T
\end{equation}



\subsubsection{核心分解}

\begin{equation}
\text{Full Gradient} =
\underbrace{g_u^l[t]}_{\text{空间}}
\times
\underbrace{\hat{a}^{l-1}[t]}_{\text{时间}}
\end{equation}



\subsection{关键突破一：Temporal Dependency 重构}

\subsubsection{BPTT中的时间依赖}

定义离散时间依赖：

\begin{equation}
\epsilon^l[t]
=
\frac{\partial u^l[t+1]}{\partial u^l[t]}
+
\frac{\partial u^l[t+1]}{\partial s^l[t]}
\frac{\partial s^l[t]}{\partial u^l[t]}
\end{equation}



\subsubsection{NDOT提出连续形式}

\begin{equation}
e(t) \triangleq \frac{\partial u(t+1)}{\partial u(t)}
\end{equation}



\subsubsection{推导（关键步骤）}

由链式法则：

\begin{equation}
\frac{d u(t+1)}{dt}
=
\frac{\partial u(t+1)}{\partial u(t)}
\frac{d u(t)}{dt}
\end{equation}

得到：

\begin{equation}
e(t) =
\frac{u'(t+1)}{u'(t)}
\end{equation}



\subsubsection{代入神经动力学}

由LIF模型：

\begin{equation}
e(t) =
\frac{u(t+1) - I(t+1)}{u(t) - I(t)}
\end{equation}



\subsubsection{离散化结果（核心公式）}

\begin{equation}
e^l[t]
=
\frac{u^l[t] - V_{th} s^l[t]}
{u^l[t-1] - V_{th} s^l[t-1]}
\end{equation}



\subsubsection{关键意义}

\begin{itemize}
\item 用神经元动力学替代BPTT链式展开
\item 不需要存储所有历史
\item 时间依赖被“隐式编码”
\end{itemize}



\subsection{关键突破二：Temporal Gradient递推}

定义：

\begin{equation}
\hat{a}^{l-1}[t]
=
s^{l-1}[t]
+
\sum_{k<t}
P^{l-1}_{k,t} \odot s^{l-1}[k]
\end{equation}

其中：

\begin{equation}
P^{l}_{k,t}
=
\prod_{i=k}^{t-1} e^l[i]
\end{equation}



\subsubsection{递推形式（核心公式）}

\begin{equation}
\hat{a}^{l-1}[t]
=
e^{l-1}[t-1] \odot \hat{a}^{l-1}[t-1]
+
s^{l-1}[t]
\end{equation}



\subsubsection{推导解释}

展开递推：

\begin{align}
\hat{a}[t]
&= s[t] + e[t-1] s[t-1] + e[t-1]e[t-2] s[t-2] + \cdots
\end{align}



\subsubsection{含义}

\begin{equation}
\hat{a}[t] = \text{带衰减的历史累积}
\end{equation}



\subsubsection{本质}

NDOT的本质是什么 

\begin{equation}
\text{BPTT的链式乘积} \rightarrow \text{递推形式}
\end{equation}



\subsection{关键突破三：完整梯度表达}

\begin{equation}
\nabla_{W^l} L
=
\sum_{t=1}^{T}
g_u^l[t] \cdot \hat{a}^{l-1}[t]^T
\end{equation}



\subsubsection{空间梯度}

\begin{equation}
g_u^l[t]
=
\left(
\frac{\partial L[t]}{\partial s^N[t]}
\prod_{j=l}^{N-1}
\frac{\partial s^{j+1}[t]}{\partial u^j[t]}
\frac{\partial s^l[t]}{\partial u^l[t]}
\right)^T
\end{equation}



\subsubsection{关键性质}

\begin{itemize}
\item 不依赖未来时间
\item 可在线计算
\end{itemize}



\subsection{关键突破四：Instantaneous Loss}

\begin{equation}
L[t]
=
(1-\alpha)\ell_{CE}(s^N[t], y)
+
\alpha \ell_{MSE}(s^N[t], y)
\end{equation}



\subsubsection{意义}

\begin{itemize}
\item 避免依赖全时间序列
\item 支持在线学习
\end{itemize}



\subsection{NDOT完整更新规则}

\begin{equation}
\Delta W^l[t]
=
g_u^l[t] \cdot \hat{a}^{l-1}[t]^T
\end{equation}



\subsection{算法本质总结}

\begin{equation}
\text{NDOT} =
\text{空间梯度（局部时间）}
+
\text{时间梯度（动力学递推）}
\end{equation}


\section{OTTT方法的完整推导与信息结构分析}

\subsection{脉冲神经元模型与时间依赖}

考虑离散时间下的脉冲神经元，其状态更新满足：

\begin{equation}
u_i[t+1] = \lambda (u_i[t] - V_{th} s_i[t]) + \sum_j W_{ij}s_j[t] + b_i
\end{equation}

\begin{equation}
s_i[t+1] = H(u_i[t+1] - V_{th})
\end{equation}

该模型体现出以下两个重要性质：

\begin{itemize}
\item 状态 $u[t]$ 依赖于过去状态 $u[t-1]$（时间递归）
\item 输出 $s[t]$ 通过非线性阈值函数生成（不可导）
\end{itemize}

因此整个系统具有强烈的时间耦合特性，这也是训练困难的根本来源。



\subsection{BPTT中的完整梯度表达}

定义总损失：

\begin{equation}
L = \sum_{t=1}^{T} L[t]
\end{equation}

根据链式法则，对权重求导：

\begin{equation}
\frac{\partial L}{\partial W_l}
=
\sum_{t=1}^{T}
\frac{\partial L}{\partial s^{l+1}[t]}
\frac{\partial s^{l+1}[t]}{\partial u^{l+1}[t]}
\frac{\partial u^{l+1}[t]}{\partial W_l}
+
\sum_{t=1}^{T}
\frac{\partial L}{\partial s^{l+1}[t]}
\frac{\partial s^{l+1}[t]}{\partial u^{l+1}[t]}
\sum_{\tau < t}
\frac{\partial u^{l+1}[t]}{\partial u^{l+1}[\tau]}
\frac{\partial u^{l+1}[\tau]}{\partial W_l}
\end{equation}

该表达式包含两部分：

\begin{itemize}
\item 第一项：当前时间步的直接梯度贡献
\item 第二项：历史时间通过状态传播影响当前输出
\end{itemize}

进一步展开时间传播项：

\begin{equation}
\frac{\partial u[t]}{\partial u[\tau]} =
\prod_{i=\tau}^{t-1}
\left(
\frac{\partial u[i+1]}{\partial u[i]}
+
\frac{\partial u[i+1]}{\partial s[i]}
\frac{\partial s[i]}{\partial u[i]}
\right)
\end{equation}

该项描述了从时间 $\tau$ 到时间 $t$ 的所有路径贡献，是时间信用分配的核心。



\subsection{时间传播结构的物理含义}

上述乘积项可以解释为：

\begin{equation}
\text{时间传播} =
\text{状态传播} + \text{脉冲触发路径}
\end{equation}

其中：

\begin{itemize}
\item $\frac{\partial u[i+1]}{\partial u[i]}$：连续状态传播（leaky dynamics）
\item $\frac{\partial u[i+1]}{\partial s[i]}$：reset路径
\end{itemize}

这意味着梯度不仅沿状态传播，还沿脉冲触发路径传播。



\subsection{OTTT的关键近似：去除reset梯度路径}

OTTT做出核心近似：

\begin{equation}
\frac{\partial u[i+1]}{\partial s[i]}
\frac{\partial s[i]}{\partial u[i]} \approx 0
\end{equation}

理由如下：

\begin{itemize}
\item $H(\cdot)$ 的导数在绝大多数点为0
\item reset路径对梯度贡献极小
\end{itemize}

因此：

\begin{equation}
\frac{\partial u[i+1]}{\partial u[i]} = \lambda
\end{equation}



\subsection{时间传播的简化形式}

于是时间传播项简化为：

\begin{equation}
\frac{\partial u[t]}{\partial u[\tau]} = \lambda^{t-\tau}
\end{equation}

将其代入梯度表达式：

\begin{equation}
\frac{\partial L}{\partial W_l}
=
\sum_{t=1}^{T}
g_u[t]
\sum_{\tau \le t}
\lambda^{t-\tau} s^l[\tau]
\end{equation}

其中：

\begin{equation}
g_u[t] = \frac{\partial L}{\partial s[t]} \cdot \frac{\partial s[t]}{\partial u[t]}
\end{equation}



\subsection{Eligibility Trace的引入}

定义：

\begin{equation}
\hat{a}[t] = \sum_{\tau \le t} \lambda^{t-\tau} s[\tau]
\end{equation}

该变量具有以下意义：

\begin{itemize}
\item 表示过去所有输入的指数加权累积
\item 是时间信息的压缩表示
\end{itemize}

为了实现在线计算，将其写为递推形式：

\begin{equation}
\hat{a}[t+1] = \lambda \hat{a}[t] + s[t+1]
\end{equation}

这一递推消除了对历史序列的存储需求。



\subsection{梯度表达的最终形式}

梯度可写为：

\begin{equation}
\frac{\partial L}{\partial W_l}
=
\sum_{t=1}^{T}
g_u[t] \hat{a}[t]^T
\end{equation}

对应单个连接：

\begin{equation}
\Delta W_{ij} = \hat{a}_i[t] \cdot \delta_j[t]
\end{equation}



\subsection{瞬时损失的引入}

传统损失依赖整个时间序列：

\begin{equation}
L = L\left(\frac{1}{T} \sum_{t} s[t]\right)
\end{equation}

OTTT改为：

\begin{equation}
L = \sum_{t=1}^{T} L[t], \quad L[t] = \frac{1}{T} L(s[t], y)
\end{equation}

该修改使得梯度仅依赖当前时间步，从而实现在线更新 \cite{xiao2022online}。



\subsection{信息结构分析}

BPTT梯度可以重写为：

\begin{equation}
\frac{\partial L}{\partial W}
=
\sum_{\tau}
s[\tau]
\sum_{t \ge \tau}
\delta[t]\lambda^{t-\tau}
\end{equation}

该表达式包含：

\begin{itemize}
\item 历史输入 $s[\tau]$
\item 当前误差 $\delta[\tau]$
\item 未来误差 $\delta[t], t>\tau$
\end{itemize}

而OTTT近似为：

\begin{equation}
\frac{\partial L}{\partial W}
=
\sum_{t}
\delta[t]\hat{a}[t]
\end{equation}

即：

\begin{equation}
\text{OTTT} \approx \text{BPTT} - \text{未来误差项}
\end{equation}



\subsection{三因子学习结构}

最终更新形式：

\begin{equation}
\Delta W_{ij} = \underbrace{\hat{a}_i[t]}_{\text{历史输入}} \cdot
\underbrace{f(u_j[t])}_{\text{神经元响应}} \cdot
\underbrace{\delta_j[t]}_{\text{误差信号}}
\end{equation}

该结构对应典型三因子学习规则：

\begin{itemize}
\item presynaptic trace
\item postsynaptic activity
\item global error
\end{itemize}



\subsection{方法本质总结}

OTTT通过以下三个步骤实现在线化：

\begin{enumerate}
\item 删除时间反向传播路径（reset路径）
\item 使用eligibility trace压缩历史信息
\item 使用瞬时损失消除未来依赖
\end{enumerate}

其本质可以写为：

\begin{equation}
g_{\text{OTTT}} = g_{\text{past}} + g_{\text{present}}
\end{equation}

相比BPTT缺失：

\begin{equation}
g_{\text{future}}
\end{equation}

这一缺失构成其性能与近似误差的根本来源。


OTTT 和 NDOT 的核心区别在于：

OTTT通过假设时间梯度为一个固定衰减因子，忽略了神经元发放带来的非线性影响，从而实现了在线训练，但代价是丢失了大量时间信息。

而 NDOT 则从神经元动力学出发，用膜电位的变化关系构造时间依赖项，将复杂的时间梯度用一个动态变量 $e[t]$ 表示，从而在保持在线性的同时，更准确地逼近真实梯度。





\bibliographystyle{IEEEtran}
\bibliography{SNNInformationReconstruction}



\end{document}
